\subsection{A finite geometric quotient detecting the completed obstruction}

\begin{statement}{A.93}{Proved}
{The old completed tuple has no locally matching section}
Let
\[
 \begin{aligned}
 U_3&=\mathbf P^1_\Q\setminus\{c_1,c_2,c_3,\infty\},
 &\Pi&=\pi_1(U_{3,\overline\Q}),\\
 H&=\ker\bigl(\Pi\longrightarrow(\Z/4)^3\bigr),
 &K&=\overline{H^2[H,H]} .
 \end{aligned}
 \tag{A.93.1}
\]
where the three geometric characters are the fourth-root
characters of $T-c_i$ and the $c_i$ are those of A.90.
Use the rational point $R$ to split the arithmetic
fundamental-group sequence.
Then $K$ is normal in the arithmetic group,
\[
 |\Pi/K|=2^{135},\qquad \Pi''\subset K ,
 \tag{A.93.2}
\]
and the following obstruction already occurs in
$(\Pi/K)\rtimes G_\Q$.

There is no section of this quotient whose three
Kummer classes modulo four are $[a_i]$ and whose
restriction at every finite prime is conjugate to
the image of the ordinary point $t_p$ constructed
in A.90, with its real restriction conjugate to
the image of the ordinary point $R$.
Consequently there is no full Selmer section of
$U_3$ with the three completed coordinates $A_i$
of A.90.
There is also no full metabelian section with
these three completed coordinates and these
ordinary local point images, including the real image.
The same obstruction factors through the lower
central quotient of class eleven:
with $\Gamma_1\Pi=\Pi$ and
$\Gamma_{r+1}\Pi=\overline{[\Pi,\Gamma_r\Pi]}$,
there is an actual arithmetic quotient map
\[
       (\Pi/\Gamma_{12}\Pi)\rtimes G_\Q
                         \longrightarrow(\Pi/K)\rtimes G_\Q .
 \tag{A.93.2a}
\]
\end{statement}
\begin{proof}
We identify the cover subgroup, the finite quotient
carrying its class, and the local pullback maps.
Let $C^\circ$ be the actual finite etale cover of
$U_3$ with equations
\[
                     T-c_i=a_iX_i^4\quad(1\le i\le3).
 \tag{A.93.3}
\]
Here $C^\circ$ removes the fibres over the three
old finite punctures and infinity; it does not
remove the fibre over the new point $c$.
It has the rational point $P_0=(R,1,1,1)$.
Its geometric fundamental group is the subgroup
$H$ of (A.93.1), and its arithmetic fundamental
group, with this point splitting, is
$H\rtimes G_\Q\subset\Pi\rtimes G_\Q$.
The geometric degree is $64$.
The free profinite group $\Pi$ has rank three,
so the subgroup theorem for finite-index
subgroups of a free profinite group gives
\[
            \operatorname{rank}(H)=1+64(3-1)=129,
       \qquad H/K=(\Z/2)^{129}.
 \tag{A.93.4}
\]
One can also read the rank from the actual curve:
its genus is $33$ and it has $64$ geometric
boundary points, giving $2\cdot33+64-1=129$.
Equivalently, the finite covering graph has
$64$ vertices and $3\cdot64$ edges, so a
spanning tree leaves exactly $129$ free
generators.  The algebraic--topological
cover comparison here is
\cite[Remark 3.3(c) and Theorem 21.3]{AgentMilneLEC}.
The subgroup $H$ is invariant under the arithmetic
action, and $K$ is characteristic in $H$.
This proves arithmetic normality and the order
in (A.93.2).
Since $\Pi'\subset H$, one has
$\Pi''\subset[H,H]\subset K$.
Thus there is the explicit arithmetic quotient map
\[
       (\Pi/\Pi'')\rtimes G_\Q
                      \longrightarrow(\Pi/K)\rtimes G_\Q .
 \tag{A.93.5}
\]

We also prove the nilpotency bound in (A.93.2a).
Put $Q=\Pi/K$, $N=H/K$, and
$A=Q/N=(\Z/4)^3$.
The subgroup $N$ is elementary abelian of exponent
two, and $Q'\subset N$ because $A$ is abelian.
Conjugation on $N$ factors through $A$, so $N$
is a module over the actual group ring
$\mathbf F_2[A]$.
If $g_1,g_2,g_3$ are the three generators of $A$,
the map $x_i\mapsto g_i-1$ gives
\[
 \mathbf F_2[A]
       =\mathbf F_2[x_1,x_2,x_3]/(x_1^4,x_2^4,x_3^4).
 \tag{A.93.5a}
\]
Indeed $(g_i-1)^4=g_i^4-1=0$ in characteristic
two, and the displayed source and target both
have the $64$ indicated monomials as a basis.
For the augmentation ideal $I=(x_1,x_2,x_3)$,
every degree-ten monomial has some exponent
at least four.  Thus $I^{10}=0$.
Commutation with $Q$ on a submodule of $N$
is precisely multiplication by $I$, because
the action is the one just specified.
Consequently
\[
       \Gamma_{2+k}Q=I^k Q'\quad(k\ge0),\qquad
       \Gamma_{12}Q=I^{10}Q'=0 .
 \tag{A.93.5b}
\]
The quotient homomorphism $\Pi\to Q$ therefore
kills $\Gamma_{12}\Pi$, proving (A.93.2a).
All these subgroups are invariant under the
given arithmetic action.

\emph{The normalized Brauer class and its exact cochain.}
The unramified class
\[
          \mathcal A=\bigl((T-c)/(R-c),2\bigr)
                           \in\operatorname{Br}(C^\circ)[2]
 \tag{A.93.6}
\]
is the restriction of the class on the smooth
projective curve $C$ constructed in A.69 and A.88.
Choose a lift in $H^2_{\rm et}(C,\mu_2)$ using
the Kummer sequence, and let
$e\in H^2_{\rm et}(C^\circ,\mu_2)$ be its restriction.
This choice also gives the unramified pullback at
any real boundary point used below.
Its restriction to $P_0$ is zero:
$\mathcal A(P_0)=0$ and
$H^2(\Q,\mu_2)\longrightarrow\operatorname{Br}(\Q)[2]$
is an isomorphism, since $\operatorname{Pic}(\Q)=0$.

The natural comparison
\[
 H^2(H\rtimes G_\Q,\mathbf F_2)
                         \xrightarrow{\ \sim\ }
                         H^2_{\rm et}(C^\circ,\mu_2)
 \tag{A.93.7}
\]
identifies $\mu_2$ with its trivial two-element
Galois module.
Here is the geometric input to this comparison.
The characteristic-zero topological curve
retracts onto a finite graph; the finite-cover
comparison just cited identifies its profinite
fundamental group with $H$.
For finite locally constant coefficients the
etale--Betti comparison is
\cite[Theorems 21.1 and 21.5]{AgentMilneLEC}.
For the coefficient $\mu_2$ used here the
geometric vanishing can also be checked
directly over $\overline\Q$:
the localization sequence and the class of
any removed geometric point map onto
$H^2(\overline C,\mu_2)=\mathbf F_2$, while
the proper-curve cohomology is zero above
degree two.  The local support groups in
\cite[Propositions 14.2 and 14.3]{AgentMilneLEC}
therefore give
$H^q_{\rm et}(C^\circ_{\overline\Q},\mu_2)=0$
for $q\ge2$.
On the group side freeness gives the same
vanishing, and in degree one the natural
map is the same cover-character map
$H\to\mathbf F_2$.
The geometric comparison is also stated directly in
\cite[Proposition 14.12 and Theorem 14.14]{AgentMilneLEC}.
The two arithmetic Hochschild--Serre spectral
sequences and this comparison give (A.93.7).
Equivalently the usual etale $K(\pi,1)$
comparison follows by applying this argument
to finite covers and finite locally constant
coefficients.

We now compute the class on the left of
(A.93.7), rather than presume that its
Hochschild--Serre edge is inflated from a
particular quotient.
Represent $e$ by a central profinite extension
of $H\rtimes G_\Q$ by $\mathbf F_2$.
Its restriction to $G_\Q$ splits continuously
because $P_0^*e=0$.
Its restriction to $H$ splits continuously
because $H$ is free profinite and hence
projective.
Choose these two splittings.
For $g\in G_\Q$ and $h\in H$, their failure
to commute according to the given action is
described by a homomorphism in $h$.
Writing $gh$ for the action of $g$ on $h$,
define $\eta(g)\in\operatorname{Hom}_{\rm cont}
(H,\mathbf F_2)$ by
\[
 s_G(g)s_H(h)s_G(g)^{-1}
       =\eta(g)(gh)\,s_H(gh),
 \tag{A.93.8}
\]
where the first factor denotes the central
element using additive $\mathbf F_2$ notation.
Associativity gives
\[
                 \eta(g_1g_2)=\eta(g_1)+g_1\eta(g_2),
       \qquad (g\varphi)(h)=\varphi(g^{-1}h).
 \tag{A.93.9}
\]
This is a continuous one-cocycle with values
in the finite group
$\operatorname{Hom}(H,\mathbf F_2)$.
The continuous section
$(h,g)\mapsto s_H(h)s_G(g)$ of the extension
has the literal two-cocycle
\[
 c_\eta\bigl((h_1,g_1),(h_2,g_2)\bigr)
                          =\eta(g_1)(g_1h_2).
 \tag{A.93.10}
\]
Its cocycle identity follows by expanding
(A.93.9); both $h_2$ and $g_2h_3$ are
evaluated by the same homomorphism.
Changing the splitting over $H$ by
$\varphi\in\operatorname{Hom}(H,\mathbf F_2)$
changes $\eta$ by its one-coboundary.
Thus $[\eta]$ is precisely the
Hochschild--Serre edge of the normalized
class, in the convention (A.93.10).
There is no sign issue over $\mathbf F_2$.

Every $\eta(g)$ kills $K=H^2[H,H]$.
Consequently (A.93.10) defines a class
\[
 \bar e\in H^2((H/K)\rtimes G_\Q,\mathbf F_2),
 \qquad
      \operatorname{inf}(\bar e)=e .
 \tag{A.93.11}
\]
This is the exact finite geometric quotient
through which the class on the cover factors.
The class is on the displayed cover subgroup
of $(\Pi/K)\rtimes G_\Q$; no extension of
that class to the whole larger group is
needed for the argument.

\emph{The modulo-four coordinate condition lifts the section.}
The finite cover (A.93.3), with its rational
point over $R$, gives the actual arithmetic
quotient
\[
        \Pi\rtimes G_\Q\longrightarrow
                     (\mu_4^3)\rtimes G_\Q .
 \tag{A.93.12}
\]
For a section, its three cocycle coordinates
in this quotient are
$\kappa_s(T-c_i)-\kappa(a_i)$ modulo four.
The assumption $[a_i]$ makes their cohomology
class zero.
Conjugating by an element of $\mu_4^3$ makes
the cocycle identically zero.
The map $\Pi/K\to\mu_4^3$ is surjective, so
this conjugation lifts to $\Pi/K$.
The conjugated section therefore lands in
the inverse image of the zero section of
(A.93.12), namely
$(H/K)\rtimes G_\Q$.
Denote this lifted section by $s_C$.
Its pullback
\[
          s_C^*\bar e\in H^2(\Q,\mu_2)
                                  =\operatorname{Br}(\Q)[2]
 \tag{A.93.13}
\]
is one global Brauer class.

\emph{The actual local points in the lifted cover.}
At a finite prime $p$, assume the original
quotient section is conjugate to the image
of the ordinary point $t_p$.
The sheets of (A.93.3) over $t_p$ form the
finite $G_{\Q_p}$-set obtained by restricting
(A.93.12) along that point section.
The condition that $s_C$ lands in the cover
subgroup is exactly that the corresponding
coset is fixed in this set.
A fixed sheet is a rational point
$Q_p\in C^\circ(\Q_p)$ above $t_p$.
Choosing its corresponding sheet path makes
the local section conjugate, by $H/K$, to
the image of that point.
This is the finite-cover stabilizer
description of the local lifting, applied
to the actual displayed cover.

Pullback of (A.93.11) along the point
section now identifies the restriction of
(A.93.13) with the evaluation
$\mathcal A(Q_p)$.
All possible sheet choices have the same
evaluation, because the function in
(A.93.6) depends only on $T$.
A.90 proves that its invariant is $1/2$
at two and zero at every other finite prime.
The prescribed real point lifts by the
same fixed-sheet argument.  Its evaluation
is zero: the second quaternion entry two
is a square in $\mathbf R$.
Thus the global
class (A.93.13) has invariant sum $1/2$.
This contradicts global Brauer reciprocity
\cite[Theorem 12.1.8]{AgentCTSkoro2019}.
It excludes the asserted locally matching
quotient section.

\emph{Full Selmer sections and the exact puncture set.}
For a full Selmer section with the completed
coordinates $A_i$, the local ordinary
points have precisely the parameters
$t_p$.
The map
$\Q_p^\times\to\widehat{\Q_p^\times}$
is injective: the valuation group
$\Z$, the finite roots of unity, and
the compact principal-unit group are
each separated by their power quotients.
Thus equality of the first completed
coordinate at two ordinary points
already gives equality of their
$T$ values.

A finite cusp at $c_i$ cannot have these
completed coordinates.
For $j\ne i$, its regular unit
$T-c_j$ specializes to $c_i-c_j$.
Kummer injectivity would give
$t_p-c_j=c_i-c_j$, hence $t_p=c_i$,
contrary to (A.90.2).
At infinity the functions
$T-c_i=z^{-1}(1-c_i z)$, $z=1/T$,
all have the same tangential Kummer
class.  Equality of the first two
ordinary coordinate classes would
give $t_p/(t_p-1)=1$, also impossible.
These are the explicit regular-unit
and tangential maps excluding all
local cusps at finite primes.
At the real place the Selmer restriction
is a real point or a real tangential
cusp.  In the first case the same
fixed-sheet argument gives a real point
of $C^\circ$ and zero evaluation.
In the cusp case the lifted local
section lies in a decomposition group
of the proper curve $C$ at a real
boundary point.
Indeed, inertia orbits on the cover's
geometric sheets index its geometric
boundary points.  The lifted section
fixes its sheet and hence that inertia
orbit, which gives a real boundary
point; the map to the proper curve
kills the corresponding inertia.
Since $\mathcal A$ extends across that
boundary, its pullback through the
tangential map equals its value at the
boundary point; that value is zero
because two is a square in $\mathbf R$.
Thus either Selmer alternative gives
zero real invariant.
Thus the full Selmer section gives
the same reciprocity contradiction,
and the full metabelian consequence
with stipulated ordinary local images
follows from (A.93.5).
The class-eleven consequence follows by
the same projection through (A.93.2a).

The old punctured line in this proof
has the four punctures
$0,1,c_3,\infty$.
The forgetful map to the tripod is
the map induced by
$U_3\hookrightarrow\mathbf P^1
\setminus\{0,1,\infty\}$; it kills
the inertia at $c_3$.
The quotient $\Pi/K$ retains its
nonzero third coordinate modulo four,
so that quotient cannot factor
through the tripod fundamental group.
Accordingly the result does not
exclude a tripod section having
only the first two completed
coordinates of A.90.
\end{proof}

The real-point matching condition in the finite-quotient
assertion is included deliberately.
The equality of a Brauer class with zero after base
change to $\mathbf R$ does not alone identify the
pullback of a chosen $\mu_2$ lift along an arbitrary
section of a finite geometric quotient: the Kummer
sequence also contains
$\operatorname{Pic}(C^\circ_\mathbf R)/2$.
The actual ordinary and tangential evaluation maps
in the proof supply the needed zero real pullback.

The exact certificate
\path{code/completed_old_section_certificate.py},
with output in
\path{code/completed_old_section_certificate_output.txt},
checks the $64$-sheet covering graph and its rank,
the literal cross-cocycle and splitting-change
identities in finite semidirect products, and
the geometric quotient order.
